\chapter{Results and Analysis}
\label{ch:resultados} % la etiqueta para referencias

% {\color{blue} En esta sección de debe incluir todos los resultados obtenidos, incluyendo figuras, gráficas y tablas. Estos resultados deben ser explicados, analizados y discutidos respecto del estado del arte de la investigación. (Máximo 20 páginas). }

In this section, we provide a comprehensive analysis of the results obtained from the pre-training and fine-tuning processes. We delve into the details of hyperameters set for the different stages of training, the collected and processed data, highlighting key metrics such as dataset size and quantity. Additionally, we present the accuracy achieved by the pre-trained transformer model on the final epoch run for each dataset..... 
The following tables show the respective hyperparameters for the pre-training and fine tuning algorithms. These hyperparameters were carefully selected based on experimentation and empirical evidence to strike a balance between model convergence, computational efficiency, and resource constraints......
\begin{table}[H]
\label{tab:Pre-Training HyperParameters}
\centering
\begin{tabular}{|l|l|l|l|l|l|}
\hline
Batch Size & LR & Epochs & Warm-up Ratio & Max Length \\ \hline
80         &    1.0x$10^{-4}$    &  15   & 0.1  & 512   \\ \hline
\end{tabular}
\caption{Pre-Training Hyperparameters}
\end{table}


\vspace{0.5cm}


\begin{table}[H]
\label{tab:Fine Tuning HyperParameters}
\centering
\begin{tabular}{|l|l|l|l|l|l|l|}
\hline
Model & Batch Size & LR & Epochs & Warm-up Ratio & Max Length \\ \hline
BETO & 8   &   1.5x$10^{-6}$    &  10   & 0.2  & 512 \\ \hline
ALBERT & 2   &   1.5x$10^{-6}$    &  10   & 0.2  & 512 \\ \hline
RoBERTa & 8   &   1.5x$10^{-6}$    &  10   & 0.2  & 512   \\ \hline
RoBERTa-CHLaw & 8   &   1.5x$10^{-6}$    &  10   & 0.2  & 512  \\ \hline
\end{tabular}
\caption{Fine Tuning Hyperparameters}
\end{table}

\begin{itemize}
   \item Batch Size: The batch size determines the number of training examples processed together in a single iteration during fine-tuning.
   \item LR: The learning rate determines the step size at which the model's parameters are updated during fine-tuning.
   \item Epochs: The number of epochs represents the total number of times the entire dataset is passed through during fine-tuning.
   \item Warm-up Ratio: The warm-up ratio refers to the initial period of training with a smaller learning rate to stabilize the fine-tuning process.
   \item Max Length: The maximum length parameter determines the maximum number of tokens allowed in an input sequence.
 \end{itemize}




\section{Results and Metrics Analysis}
In this section, we delve into the outcomes of the conducted experiments and analyze the performance of various models. We begin by exploring the data that underpins both the pre-training and fine-tuning stages. Furthermore, we present an assessment of the pre-trained transformer models' efficacy, highlighting their performance metrics. Subsequently, we provide a comprehensive analysis of the fine-tuning process for each of the selected OAECAs.
\subsection{Pre-Training Transformer Results/Metrics}

The amount of computational resources available permitted pre-training on five splits at a time, each one running for fifteen epochs. The following figures show the last training epoch of the last run split of the respective data set.

% \vspace{0.5cm}
% \begin{figure}[H]
% \label{fig:Last epoch ParaCrawl Dataset}
% \centering
% \includegraphics[scale=1]{images/Paracrawl.png}
% \caption{Last epoch metrics ParaCrawl Data set}
% \end{figure}

\begin{table}[H]
\label{table:LastepochParaCrawlDataset}
\centering
\begin{tabular}{|l|l|l|l|l|l|l|l}
\hline
Accuracy & Train Loss & Eval Loss & LR & Train Time & Eval Time \\ \hline
0.598 & 2.137 & 2.167 &   2.72x$10^{-9}$ &  25:30  & 01:53   \\ \hline
\end{tabular}
\caption{Last epoch metrics ParaCrawl Data set}
\end{table}

